% Pregunta 399 (2011-octubre-OEP2011, n.º 20): consumo de estado estacionario c* = f(k*) - (n+delta)k* con f(k)=1,6*sqrt(k) y n+delta=0,40; máximo en k_oro=4 (regla de oro); a la derecha, más capital da menos consumo. Autor: Víctor Gutiérrez Marcos
\documentclass[tikz,border=6pt]{standalone}
\input{tikz-tcee}
\begin{document}
\begin{tikzpicture}[tcee,x=0.5cm,y=2.4cm]
  \ejes{$k^*$}{$c^*$}{12}{2}
  \fill[tceerojo,opacity=0.08] (4,0) rectangle (11.5,1.85);
  \draw[curva1,domain=0:11.5,samples=100] plot (\x,{1.6*sqrt(\x)-0.4*\x}) node[etiqueta,right] {$c^*$};
  \draw[guia] (4,1.6) -- (4,0) node[below] {$k_{\text{oro}}$};
  \draw[guia] (4,1.6) -- (0,1.6) node[left] {$c_{\text{oro}}$};
  \node[punto] at (4,1.6) {};
  \node[etiqueta,anchor=south,font=\footnotesize] at (4,1.63) {$f'(k_{\text{oro}})=n+\delta$};
  \draw[flecha,tceeverde] (1.2,0.72) -- (2.6,1.18);
  \node[etiqueta,tceeverde,font=\footnotesize,align=center] at (2.4,0.5) {$k\uparrow\Rightarrow c\uparrow$};
  \draw[flecha,tceerojo] (6.4,1.66) -- (9.6,1.33);
  \node[etiqueta,tceerojo,font=\footnotesize,align=center] at (8.4,0.62) {$k\uparrow\Rightarrow c\downarrow$\\(ineficiencia\\dinámica)};
\end{tikzpicture}
\end{document}
